\clearpage
\begin{summarybox}
\textbf{\textcolor{CSummary}{考场 5 秒快记}}\par
\[
\boxed{y^2=\pm2px\ (p>0),\quad
L_{\text{焦点弦}}=\frac{2p}{\sin^2\theta}\ge2p.}
\]
\textbf{通径：}$\theta=90^\circ$，$L=2p$（最短）。\par\par\medskip
\textbf{\textcolor{CSummary}{忘了公式？从“最熟悉”的三步救回来}}\par
\begin{enumerate}[leftmargin=*,itemsep=.25em]
\item 焦点弦两端的\textbf{焦半径相加}就是整条弦：$L=AF+BF$。
\item 抛物线定义：\textbf{到焦点 = 到准线}；一端有 $r=p+r\cos\theta$，另一端有 $s=p-s\cos\theta$。
\item \textbf{一加一化：}
$L=\frac p{1-c}+\frac p{1+c}=\frac{2p}{1-c^2}=\frac{2p}{\sin^2\theta}$，$c=\cos\theta$。
\end{enumerate}
\textbf{\textcolor{CSummary}{如果投影不熟：最低门槛的代数备用}}\par
$L=x_A+x_B+p$；设 $x=p/2+y\cot\theta$；由抛物线求出
$y_A+y_B=2p\cot\theta$，故 $L=2p(1+\cot^2\theta)$。
\par\medskip
\textbf{\textcolor{CSummary}{题型识别与检验}}\par
\begin{itemize}[leftmargin=*,itemsep=.15em]
\item 给角度：直接代；给斜率：$L=2p(1+k^2)/k^2$（有限 $k\ne0$）；给弦长：反求 $\sin^2\theta=2p/L$，通常有两个互补倾角。
\item 左右开口统一；上下开口（角仍从 $x$ 轴量）改为 $\cos^2\theta$。
\item \textbf{结果低于 $2p$ 必错；竖直焦点弦必须是 $2p$。}
\end{itemize}
\end{summarybox}
